%----------------------------------------------------------------------------
\chapter{\bevezetes}

A modern tudomány egyik legnagyobb kihívása a komplex biológiai rendszerek dinamikájának megértése és előrejelzése. Az ökoszisztémák vizsgálata a valóságban gyakran akadályokba ütközik: a folyamatok időigényesek, a beavatkozások pedig visszafordíthatatlanok vagy etikailag megkérdőjelezhetőek lehetnek. A számítógépes szimulációk ezzel szemben biztonságos, költséghatékony és reprodukálható környezetet biztosítanak a populációgenetikai és ökológiai elméletek teszteléséhez.

A diplomamunka alapvető motivációja egy olyan digitális laboratórium létrehozása volt, amely képes vizuálisan és adat szintjén is reprezentálni a természetes szelekciót és a fajok közötti interakciókat. A projekt célja nem csupán egy látványos vizualizáció elkészítése, hanem egy olyan robusztus szoftverarchitektúra kialakítása, amely alkalmas lehet tudományos igényű adatgyűjtésre és elemzésre.

A projektem gyökerei a BSc képzés idejére nyúlnak vissza, ahol az első ágensalapú modellek implementálása megtörtént. A koncepció kidolgozásakor inspirációként szolgált Sebastian Lague ismert ökoszisztéma-szimulációs projektje, amely rávilágított a Unity3D motorban rejlő lehetőségekre a biológiai modellezés terén.

Míg a kezdeti inspiráció a vizuális reprezentációra és az alapvető túlélési logikára fókuszált, a saját fejlesztésem hamar túllépett ezen a kereten. Az eddigi megoldásokhoz képest a hangsúly eltolódott a szoftvertechnológiai precizitás, a skálázhatóság és a professzionális adatkezelés irányába. A korábbi monolitikus struktúrát felváltotta egy moduláris felépítés, amely lehetővé tette a komplexebb, fajspecifikus viselkedésformák és az eseményvezérelt működés integrálását.

A dolgozatban bemutatott rendszer egy Unity3D alapú, eseményvezérelt ökoszisztéma-szimuláció. A keretrendszer legfőbb újdonsága a populációgenetikai folyamatok mély integrációja: az egyedek DNS-alapú öröklődési mechanizmusokkal rendelkeznek, amelyek meghatározzák fizikai jellemzőiket és viselkedési stratégiáikat.

A szimuláció során keletkező hatalmas mennyiségű adat hatékony kezelésére és valós idejű monitorozására egy InfluxDB adatbázisra épülő naplózó rendszert fejlesztettem ki. Az események továbbítása HTTP-alapú API hívásokon keresztül történik, az adatok vizuális megjelenítését és elemzését pedig egy Grafana dashboard biztosítja. Ez a felépítés lehetővé teszi, hogy a populációk változásait, a mutációk hatásait és a ragadozó-préda dinamikát ne csak a Unity környezetben, hanem külső statisztikai eszközökkel validált, interaktív grafikonok formájában is elemezhessük.

A dolgozat szerkezete a következőképpen alakul: A \textbf{\ref{chap:elemzes}. fejezet} a feladatkiírás részletes elemzésével és a követelmények specifikációjával foglalkozik, majd a \textbf{\ref{chap:irodalom}. fejezet} mutatja be az elméleti hátteret, a hasonló megoldásokat (különös tekintettel Sebastian Lague munkásságára), valamint az olyan biológiai alapmodelleket, mint a Lotka–Volterra egyenletek. A \textbf{\ref{chap:tervezes}. fejezet} tartalmazza a rendszertervet, a szoftverarchitektúra döntési pontjait és az eseményvezérelt működés logikáját. A megvalósítás technikai lépéseit a Unity környezetben, kitérve a faj-specifikus viselkedések programozására, a \textbf{\ref{chap:implementacio}. fejezet} részletezi. Végezetül a \textbf{\ref{chap:ertekeles}. fejezet} a szimuláció futtatása során nyert adatok statisztikai elemzését, a rendszer validálását és a továbbfejlesztési lehetőségek bemutatását tartalmazza.